\documentclass[11pt]{article}

\usepackage[margin=1in]{geometry}
\usepackage{amsmath,amssymb,amsthm}
\usepackage{enumitem}
\usepackage[colorlinks=true,linkcolor=blue,urlcolor=blue]{hyperref}

\newcommand{\MSO}{\ensuremath{\mathrm{MSO}}}
\newcommand{\MSOtwo}{\ensuremath{\mathrm{MSO}_2}}

\title{Toward a Self-Contained Tree-Automaton Theorem\\
  \large Components needed for \texttt{lecture\_note\_expanded.tex}}
\author{}
\date{}

\begin{document}
\maketitle

The main note \texttt{lecture\_note\_expanded.tex} isolates the effective
\MSO{}--tree-automaton correspondence (Thatcher--Wright, Doner) as a cited
external theorem.  At present neither \textsc{mathlib} nor \textsc{cslib} provides
a tree-automata API, so a fully self-contained (and eventually formalised)
development must build this theorem from scratch.  Only the direction used by
Courcelle's algorithm is needed: from an \MSO{} sentence over the tree vocabulary
one computes a \emph{deterministic} bottom-up tree automaton recognising its
models, runnable in linear time; the converse (automaton \(\Rightarrow\)
\MSO{}-definable) is not required.  The components below should eventually be added
to \texttt{lecture\_note\_expanded.tex}.

\section{Objects and definitions}
\begin{itemize}[leftmargin=*]
  \item Ranked trees, in particular full binary trees, and the encoding of a
    \(\Sigma\)-tree of outdegree at most two as a full binary ranked tree by
    padding each absent child with a null symbol \(\bot\).
  \item Nondeterministic and deterministic bottom-up finite tree automata:
    states, transition function, run, and accepted language.
  \item The \emph{extended alphabet} \(\Sigma\times\{0,1\}^k\), encoding an
    assignment of \(k\) free variables as extra label tracks; a first-order
    variable is encoded by a track required to carry exactly one \(1\).
\end{itemize}

\section{Base automata (atomic formulas)}
Automata recognising each atomic constraint over the extended alphabet:
\(\operatorname{child}_1(x,y)\), \(\operatorname{child}_2(x,y)\), membership
\(x\in X\), equality \(x=y\), the label predicates \(P_a\), and the
``singleton'' constraint that a first-order track carries exactly one \(1\).

\section{Closure operations (the technical core)}
\begin{itemize}[leftmargin=*]
  \item Product construction for \(\wedge\) and \(\vee\).
  \item Projection (erasing one track) for \(\exists\), which introduces
    nondeterminism.
  \item \textbf{Determinisation} by the bottom-up subset construction, with the
    finiteness of the powerset state space.  This is the crux, needed both for
    complementation and for the final deterministic automaton.
  \item Complementation for \(\neg\), via determinisation followed by swapping
    the accepting set.
\end{itemize}

\section{Compilation theorem}
By induction on the structure of an \MSO{} formula \(\varphi\) with \(k\) free
variables, construct an automaton \(A_\varphi\) over \(\Sigma\times\{0,1\}^k\)
accepting exactly the pairs (tree, assignment) satisfying \(\varphi\), and prove
correctness by the same induction from the base automata and the closure
operations.  Determinising \(A_\varphi\) yields the required deterministic
automaton.

\section{Faithfulness of the ranked encoding}
Show that under the \(\bot\)-padding encoding the relations
\(\operatorname{child}_1,\operatorname{child}_2\) and the label predicates are
first-order definable, so satisfaction is preserved between a \(\Sigma\)-tree and
its ranked encoding.

\section{Execution and complexity}
A computable deterministic bottom-up evaluation function together with a proof
that it decides acceptance.  The linear-time bound on the run is a metatheoretic
claim about this evaluation; making it precise rather than merely ``computable''
requires a cost model and is the least mechanisable part.

\section{Scope remark}
The development is sizeable---essentially the tree version of the
B\"uchi--Elgot--Trakhtenbrot theorem.  Until it is carried out, the cited theorem
should be regarded as the single substantial external dependency of the note.

\end{document}
